% =====================================================================
\chapter{Discusión y conclusiones}\label{cap:conclusiones}
% =====================================================================

\section{Conclusiones}

\begin{enumerate}
  \item El modelo de \textcite{chen2021} se deduce por completo de la teoría escalar de la
  difracción: la función de transferencia del espacio libre sale de la ecuación de Helmholtz
  (\cref{sec:espectro-angular}), y el semiángulo $\varphi_{\max}=\arcsin[\lambda/(2d_f)]$ es el
  ángulo asociado a la frecuencia de Nyquist de la malla de neuronas, no al primer cero de
  difracción de una neurona (\cref{sec:neurona}).
  \item El sistema opera en campo cercano ($N_F\approx126$) y fuera de la validez de la
  aproximación de Fresnel para las componentes más inclinadas, por lo que es necesario usar
  la función de transferencia exacta (\cref{sec:fresnel}).
  \item El gradiente del entrenamiento es el término de interferencia entre la onda que viaja
  hacia adelante y una onda de error que viaja hacia atrás con la función de transferencia
  conjugada; la deducción coincide con la diferenciación automática a precisión de máquina
  (\cref{sec:gradiente}).
  \item El criterio de conectividad del artículo es idéntico a la condición de muestreo del
  espectro angular de banda limitada (\cref{prop:identidad}) y, en el régimen paraxial,
  equivale a una condición sobre el número de Fresnel de una neurona (\cref{prop:fresnel}).
  \item La simulación reproduce las precisiones del artículo (\cref{tab:fig6}), confirma la
  hipótesis de conectividad acumulada mediante el barrido de capas $L\in\{1,3,5,8\}$
  (\cref{sec:barrido}), detecta dos inconsistencias en su Fig.~6 (\cref{sec:inconsistencias}) y
  reproduce el sistema completo de $1000\times1000$ neuronas y 11 clases en GPU con
  \SI{90.63}{\percent} de precisión en prueba (\cref{sec:sistema-completo}).
  \item La alineación lateral de las capas —una sola neurona de desplazamiento baja la precisión a
  $\approx\SI{84}{\percent}$ en el caso~2 y $\approx\SI{81}{\percent}$ en el sistema completo— y
  la separación axial $D$, cuya tolerancia relativa escala como $1/\sqrt N$ (\SI{0.10}{\percent}
  de $D$ en el sistema completo), son las tolerancias de fabricación más exigentes y explican la
  brecha entre simulación y experimento (\cref{sec:fabricacion,sec:sistema-completo}).
\end{enumerate}

\section{Limitaciones}

\begin{itemize}
  \item \textbf{Aproximación escalar.} Es razonable para $d_f\approx6\lambda$, pero no para el
  caso~1 ($d_f=0.53\lambda$), donde los efectos vectoriales y de acoplamiento entre neuronas
  vecinas pueden ser importantes \parencite[sec.~3.2]{goodman2005}.
  \item \textbf{Capas delgadas.} Cada capa se modela como una máscara infinitamente delgada; un
  relieve de \SI{0.7}{\micro m} junto a neuronas de \SI{4}{\micro m} no es estrictamente delgado.
  \item \textbf{Elementos del artículo no especificados.} La forma de los detectores, la
  normalización de la pérdida, la codificación de la entrada y el conjunto de dígitos
  alterados se tuvieron que elegir (\cref{tab:decisiones}), y pueden explicar parte de las
  diferencias.
  \item \textbf{Sin experimento.} Todos los resultados son numéricos.
\end{itemize}

\section{Trabajo futuro}

\begin{itemize}
  \item \textbf{Entrenamiento robusto ante imperfecciones:} incorporar aleatoriamente durante el
  entrenamiento desplazamientos laterales, variaciones en la distancia $D$ y cuantización de
  fase (\emph{error-aware training}) para cerrar la brecha entre simulación y experimento.
  \item \textbf{Modelado electromagnético vectorial:} emplear métodos FDTD o RCWA en arreglos
  con $d_f\lesssim\lambda$, donde el acoplamiento cercano y el relieve físico finito dejan de
  ser despreciables.
  \item \textbf{Fabricación experimental e iluminación policromática:} evaluar el desempeño
  de las máscaras cuantizadas en sílice fundida ante desalineaciones multigrado y múltiples
  longitudes de onda.
\end{itemize}

\section{Valoración final}

Este proyecto muestra que un problema moderno de cómputo óptico se explica por completo con las
herramientas fundamentales de la óptica ondulatoria: la ecuación de Helmholtz, el espectro
angular de ondas planas, el teorema de muestreo de Shannon--Nyquist y la escala fijada por el
número de Fresnel. Lejos de ser una caja negra numérica, cada regla de diseño y cada límite de
resolución de una D\textsuperscript{2}NN tienen una raíz física transparente: desde la
identidad exacta entre el criterio de conectividad y el límite de banda del propagador, hasta
el papel de la interferencia holográfica en el cálculo de su gradiente. En cuanto a la promesa
de la computación óptica, las D\textsuperscript{2}NN resultan especialmente atractivas para
tareas de procesamiento visual y filtrado pasivo de baja latencia —donde la luz clasifica o
transforma un frente de onda sin consumo energético adicional durante la inferencia—, aunque
su escalamiento práctico frente a arquitecturas electrónicas seguirá condicionado por la
linealidad de la propagación libre, la tolerancia de alineación sub-micrométrica y el cuello de
botella optoelectrónico en las etapas de entrada y detección.
